\documentclass[11pt,letterpaper]{article}
\usepackage[T1]{fontenc}
\usepackage[utf8]{inputenc}
\usepackage{lmodern}
\usepackage[margin=0.86in]{geometry}
\usepackage{microtype}
\usepackage{booktabs,longtable,tabularx,array}
\usepackage{enumitem}
\usepackage{xcolor}
\usepackage{hyperref}
\usepackage{xurl}
\definecolor{DocBlue}{HTML}{16435F}
\hypersetup{colorlinks=true,urlcolor=DocBlue,linkcolor=DocBlue}
\setlength{\parindent}{0pt}
\setlength{\parskip}{0.54em}
\setlist[itemize]{leftmargin=*,itemsep=3pt}
\begin{document}
\begin{center}
{\LARGE\bfseries Delivery Backlog and Engineering Acceptance Criteria}\par
\vspace{5pt}{\large Epics, story points and traceable implementation gates}\par
\vspace{9pt}{\small Engineering reference | Proposed design | Not a deployment or approval record}\par
\end{center}
\tableofcontents
\clearpage
\section{Sizing and readiness}
Story points are planning estimates (1,2,3,5,8) rather than labor hours. Definition of Ready: named owner, repository boundary, input fixtures, source citations, acceptance and security criteria, dependencies. Definition of Done: reviewed code/docs, tests pass, evidence attached, controls unchanged.


\section{Epic A — inventory and blockers}
\begin{itemize}
\item APPSEC-A01 | 5 points | P0 | As a platform engineer, I want complete read-only repository inventory so that security assessments have a documented denominator. Given partial API access, when collection ends, then denied and missing repos are distinct from absent alerts.
\item APPSEC-A02 | 5 points | P0 | As a developer, I want an accurate merge-block reason so that I can resolve the correct control. Given a required AppSec review and a passing CodeQL job, when classified, then REQUIRED\_REVIEW is returned without suggesting dismissal.
\item APPSEC-A03 | 3 points | P1 | As an AppSec engineer, I want monorepo owner reconciliation so that alerts reach accountable teams. Given ambiguous components, when joined, then uncertainty is retained and escalated.
\end{itemize}

\section{Epic B — policy and evidence}
\begin{itemize}
\item APPSEC-B01 | 8 points | P0 | As a security policy owner, I want deterministic versioned decisions so that approval is reproducible. Given no scan evidence, when evaluated, then SCAN\_INCOMPLETE takes precedence.
\item APPSEC-B02 | 5 points | P0 | As a security reviewer, I want independently verified exception matching so that false positives are not rubber-stamped. Given mismatched fingerprint or expired approval, when evaluated, then the request is escalated.
\item APPSEC-B03 | 3 points | P1 | As an auditor, I want input hash and policy version evidence so that decisions are traceable. Given two identical fixture inputs, when evaluated, then the normalized decision digest is identical.
\end{itemize}

\section{Epic C — integration and controls}
\begin{itemize}
\item APPSEC-C01 | 5 points | P0 | As a platform engineer, I want distinct GitHub App permission profiles so that discovery cannot dismiss alerts. Given a read-only token, when a mutation is attempted in a simulated test, then no write occurs.
\item APPSEC-C02 | 8 points | P1 | As an authorized reviewer, I want delegated dismissal simulation so that request-to-approval boundaries are validated. Given self-approval, when evaluated, then authorization is rejected.
\item APPSEC-C03 | 5 points | P1 | As an AppSec lead, I want recovery and audit procedures so that incorrect decisions can be contained. Given a revoked exception, when reconciled, then reopen/escalation is recorded.
\end{itemize}

\section{Epic D — adoption and measurement}
\begin{itemize}
\item APPSEC-D01 | 3 points | P1 | As a delivery lead, I want a nonproduction pilot to compare wait times. Given representative baseline/pilot groups, when evaluated, then denominators, medians and P95 are documented.
\item APPSEC-D02 | 3 points | P2 | As an engineer, I want accurate printable assessment pairs so that teams can gather independent evidence. Given a Blank, when rendered, then no fake approval or sample responses are printed.
\end{itemize}

\section{Dependencies and sprint proposal}
Sprint 1: A01/A02 and baseline; Sprint 2: B01/B02 and risk model; Sprint 3: C01/C02 and offline integration simulation; Sprint 4: D01/D02 and documentation quality. AppSec reviews and external deployment authorization are explicit dependencies, never assumed completed.


\section*{Document status and provenance}
This document is a proposed reference design. Actual GitHub feature availability, organization policies, approvals, and runtime performance must be independently verified before adoption. The architecture inventory is not proof of vulnerability exploitability or correctness of a dismissal.
\end{document}
